% !TeX root = ../main.tex
\section{Methodology}\label{sec:method}
\rubric{Baseline, correctness, reproducibility}{15}{Reference validation, controlled timing,
reproducible scripts and a documented environment (with \cref{sec:correctness,sec:baseline}).}

\subsection{Platform}
\newcommand{\acspinning}{Each run is pinned to one physical P-core
(\code{--cpu 2}, \code{sched_setaffinity}); its SMT sibling is left idle and no other
pinning/isolation is used.}
%% Measurement environment table (common requirement: measurement discipline).
%% Values come from common/system_info.tex. Before \input-ing this file a project may define
%%   \acspinning   how runs were pinned/isolated
%%   \acsenvextra  extra rows, e.g.  \newcommand{\acsenvextra}{SSD & model, firmware, PCIe link \\}
\providecommand{\acspinning}{\todo{e.g.\ \code{taskset -c 2} on a P-core; isolcpus if used}}
\providecommand{\acsenvextra}{}
\begin{table}[tbp]
  \centering\small
  \caption{Measurement environment.}\label{tab:environment}
  \begin{tabularx}{\linewidth}{@{}>{\raggedright\arraybackslash}p{0.2\linewidth}X@{}}
    \toprule
    CPU               & \sysCPU{} (\sysArch) \\
    SIMD ISA          & \sysISA \\
    Topology          & \sysCores; NUMA: \sysNUMA \\
    SMT               & \sysSMT \\
    Caches            & \sysCaches \\
    Memory            & \sysMemory; DIMMs: \sysDIMM \\
    Pages             & base \sysPageSize; THP: \sysHugePages \\
    OS / kernel       & \sysOS; \sysKernel \\
    Frequency policy  & governor: \sysGovernor; driver: \sysDriver; turbo: \sysTurbo; EPP: \sysEPP \\
    Pinning / isolation & \acspinning; \sysIsolation \\
    Compilers         & \sysGCC; \sysClang \\
    Tools             & \sysPerf{} (perf\_event\_paranoid \sysPerfParanoid); \sysFio; MLC: \sysMLC \\
    GPU (optional)    & \sysGPU; CUDA \sysCUDA \\
    \acsenvextra
    Collected         & \sysCollected \\
    \bottomrule
  \end{tabularx}
\end{table}


This CPU (Comet Lake, client Skylake core) implements SSE/AVX/AVX2 and FMA but \emph{not}
AVX-512, so the widest usable vector register is \qty{256}{bit} (\code{ymm}): \num{8} lanes of
float32 or \num{4} of float64. With \code{-march=native} GCC targets AVX2+FMA and prefers
\qty{256}{bit} vectors; the default (generic \code{x86-64}) target only guarantees SSE2, so the
plain auto-vectorized build emits \qty{128}{bit} (\code{xmm}) code. That 256-vs-128-bit gap is the
main reason we keep \code{autovec} and \code{native} as separate variants (\cref{tab:variants}).

\subsection{Kernels}
\begin{table}[tbp]
  \centering\small
  \caption{Kernels and floating-point operations (FLOPs) per processed element.}\label{tab:kernels}
  \begin{tabular}{@{}lll@{}}
    \toprule
    Kernel & Definition & FLOPs / element \\
    \midrule
    SAXPY               & $y_i \leftarrow a\,x_i + y_i$                         & 2 \\
    Dot product         & $s \leftarrow \sum_i x_i\, y_i$                        & 2 \\
    Elementwise multiply & $z_i \leftarrow x_i\, y_i$                            & 1 \\
    3-point stencil     & $y_i \leftarrow a\,x_{i-1} + b\,x_i + c\,x_{i+1}$      & 5 \\
    \bottomrule
  \end{tabular}
\end{table}
We count a multiply and an add as one FLOP each, so a fused multiply--add is two FLOPs; SAXPY and
the dot product are therefore 2~FLOP/element (one FMA), elementwise multiply is 1, and the stencil
is 5 (three multiplies, two adds). All four kernels from the project brief are implemented and
validated; SAXPY is used as the representative kernel for the locality, alignment and stride
studies because it has the simplest unit-stride access and a well-known arithmetic intensity.

\subsection{Build variants}
\begin{table}[tbp]
  \centering\small
  \caption{Build variants. Only the kernel translation unit (\code{src/kernels.c}) differs; the
  timing harness (\code{main.c}, \code{common.c}) and the reference (\code{reference.c}) are
  compiled identically across variants.}\label{tab:variants}
  \begin{tabularx}{\linewidth}{@{}lX@{}}
    \toprule
    Variant & Kernel compiler flags (from \code{simd_bench --info}) \\
    \midrule
    scalar    & \code{-O3 -fno-tree-vectorize -fno-tree-slp-vectorize} \\
    autovec   & \code{-O3} (generic \code{x86-64} $\Rightarrow$ SSE2, \qty{128}{bit}) \\
    native    & \code{-O3 -march=native} (AVX2+FMA, \qty{256}{bit}) \\
    native512 & not available on this CPU (no AVX-512; \qty{256}{bit} is the widest ISA) \\
    \bottomrule
  \end{tabularx}
\end{table}
The primary compiler is GCC \sysGCC{}. The reference implementation additionally uses
\code{-O2 -fno-tree-vectorize -ffp-contract=off} so it evaluates each expression in strict,
non-fused order and stays a trustworthy oracle. We deliberately do \emph{not} pass
\code{-ffast-math} to any variant: this keeps float results reproducible but, as
\cref{sec:vectorization,sec:negative} show, it also prevents the compiler from reassociating the
dot-product reduction---an effect we wanted to measure rather than hide.

\subsection{Timing and statistics}
\begin{itemize}
  \item Timer: \code{clock_gettime(CLOCK_MONOTONIC_RAW)} (resolution \qty{1}{\nano\second} on this
        host, from \code{simd_bench --info}). Each repetition times a calibrated batch of kernel
        invocations lasting at least \qty{25}{\milli\second}, i.e.\ roughly $2.5\times10^{7}$ timer
        ticks, so timer resolution and call overhead are negligible.
  \item Warm-up: 3 untimed invocations per configuration; repetitions: 12 per configuration.
  \item Reported statistic: median; error bars show the interquartile range (IQR) across the 12
        repetitions. The median resists the occasional scheduler/turbo outlier better than the mean.
  \item Pinning, SMT, governor and turbo: \cref{tab:environment}. Turbo is left enabled and the
        governor is \sysGovernor{}; the resulting frequency variation is a documented limitation
        (\cref{sec:anomalies}).
  \item Every kernel result is written through a \code{volatile} sink (\code{acs_consume_*}) and the
        kernels live in their own translation unit, so no loop is eliminated as dead code.
  \item Configuration order: the sweep is run as nested loops over size and variant
        (\code{scripts/run.sh}); it is not shuffled, which we note as a limitation.
\end{itemize}

\subsection{Data initialization and working-set sizes}
Inputs are filled with a deterministic linear-congruential sequence in $[-0.5, 0.5)$ (fixed seed),
which avoids denormals, NaNs and overflow; SAXPY's in-place update grows $y$ only linearly in the
repetition count, so it never overflows float32. The gather pattern uses a Fisher--Yates
permutation of $[0,N)$ so every element is touched exactly once with no locality.

Each kernel streams two or three arrays, so the bytes touched are (arrays)$\times$(element
size)$\times N$. For float32 SAXPY (two arrays, \qty{8}{\byte}/element) the working set crosses
this machine's \qty{32}{\kibi\byte} L1d near $N\approx\num{4096}$, the \qty{256}{\kibi\byte} L2 near
$N\approx\num{32768}$, and the \qty{20}{\mebi\byte} L3 near $N\approx\num{2.6e6}$; beyond that it is
DRAM-resident. These boundaries are annotated on the locality plot (\cref{sec:locality}).
